\documentclass[11pt,landscape,letterpaper]{article}

\usepackage[margin=1.05cm]{geometry}
\usepackage[utf8]{inputenc}
\usepackage[T1]{fontenc}
\usepackage[spanish]{babel}
\usepackage{amsmath, amssymb}
\usepackage{adjustbox}
\usepackage{tikz}
\usetikzlibrary{arrows.meta, positioning, shapes.geometric, shapes.misc, calc}
\pgfdeclarelayer{bg}
\pgfsetlayers{bg,main}

\pagestyle{empty}

% ------------------------
% COLORES
% ------------------------
\definecolor{azul}{RGB}{31,78,121}
\definecolor{azulclaro}{RGB}{221,235,247}
\definecolor{verde}{RGB}{226,239,218}
\definecolor{amarillo}{RGB}{255,242,204}
\definecolor{naranja}{RGB}{252,228,214}
\definecolor{rojo}{RGB}{244,204,204}
\definecolor{gris}{RGB}{242,242,242}
\definecolor{morado}{RGB}{234,209,220}

% ------------------------
% ESTILOS TIKZ
% ------------------------
\tikzset{
    startstop/.style={
        rounded rectangle,
        draw=azul,
        fill=azulclaro,
        very thick,
        text width=5.6cm,
        minimum height=1.0cm,
        align=center,
        font=\bfseries\small
    },
    process/.style={
        rectangle,
        rounded corners=2pt,
        draw=black!70,
        fill=gris,
        thick,
        text width=6.0cm,
        minimum height=1.05cm,
        align=center,
        font=\small
    },
    subproc/.style={
        rectangle,
        rounded corners=2pt,
        draw=black!65,
        fill=white,
        thick,
        text width=7.0cm,
        minimum height=1.2cm,
        align=center,
        font=\tiny
    },
    decision/.style={
        diamond,
        aspect=2.1,
        draw=black!70,
        fill=amarillo,
        thick,
        text width=4.2cm,
        align=center,
        inner sep=1pt,
        font=\small\bfseries
    },
    output/.style={
        rectangle,
        rounded corners=2pt,
        draw=black!70,
        fill=verde,
        thick,
        text width=6.0cm,
        minimum height=1.0cm,
        align=center,
        font=\small
    },
    dep/.style={
        rectangle,
        rounded corners=2pt,
        draw=black!70,
        fill=morado,
        thick,
        text width=7.0cm,
        minimum height=1.15cm,
        align=center,
        font=\tiny
    },
    risk/.style={
        rectangle,
        rounded corners=2pt,
        draw=black!75,
        fill=naranja,
        thick,
        text width=6.0cm,
        minimum height=1.0cm,
        align=center,
        font=\small
    },
    note/.style={
        rectangle,
        rounded corners=2pt,
        draw=azul,
        fill=azul!7,
        thick,
        text width=17.3cm,
        align=center,
        font=\scriptsize
    },
    arrow/.style={
        -{Latex[length=2.5mm]},
        thick,
        draw=black!75
    },
    dashedarrow/.style={
        -{Latex[length=2.5mm]},
        thick,
        dashed,
        draw=black!65
    }
}

\begin{document}

\begin{center}
{\Large\bfseries Diagrama de procedimientos operativos: Portafolio de bonos con CreditMetrics}\\[1mm]
{\small Portafolio de tres bonos globales con calificaciones A-, BBB- y B-; vencimiento de 5 años; valor nominal USD \$100}
\end{center}

\vspace*{\fill}

\noindent\makebox[\textwidth][c]{%
\begin{adjustbox}{center,max totalsize={\textwidth}{0.91\textheight}}
\begin{tikzpicture}[node distance=0.62cm and 1.00cm]

% =========================================================
% BLOQUE SUPERIOR: DATOS E INSUMOS
% =========================================================

\node[startstop] (inicio) {
Inicio\\
Se tiene un portafolio compuesto de tres bonos globales
};

\node[process, below=of inicio] (datos) {
Identificar características de los bonos\\
A-, BBB- y B-; vencimiento 5 años; valor nominal USD \$100; cupón anual 7.5\%, 9\% y 11\%
};

\node[subproc, below=of datos] (insumos) {
\textbf{Obtener insumos del modelo CreditMetrics}\\[-1mm]
1. Obtener la matriz de transición a 1 año.\\
2. Obtener la curva de los bonos libre de riesgo.\\
3. Obtener la curva de los bonos por calidad crediticia:
\[
\begin{aligned}
\text{tasa por calidad}
&=\text{tasa libre de riesgo}\\
&\quad+\text{sobretasa}
\end{aligned}
\]
4. Obtener la tasa de recuperación en caso de incumplimiento:
\[
\begin{aligned}
RR_{A^-}&=45\%\\
RR_{BBB^-}&=35\%\\
RR_{B^-}&=25\%
\end{aligned}
\]
};

\node[decision, below=of insumos, yshift=-0.1cm] (tipo) {
Supuesto sobre probabilidades de incumplimiento
};

% =========================================================
% RAMA IZQUIERDA: INDEPENDENCIA
% =========================================================

\node[process, below left=1.05cm and 5.1cm of tipo] (ind0) {
Caso I\\
Migraciones independientes
};

\node[subproc, below=of ind0] (ind1) {
\textbf{Valuar los bonos en cada estado de la naturaleza}\\[-1mm]
Para cada bono \(i\in\{A^-,BBB^-,B^-\}\) y cada estado \(s\):
\[
V_{i,s}
=
\sum_{\tau=1}^{4}
\frac{C_i}{(1+y_s)^\tau}
+
\frac{100+C_i}{(1+y_s)^4}
\]
donde al horizonte de 1 año quedan 4 años al vencimiento.\\
Si \(s=D\):
\[
V_{i,D}=RR_i\times 100.
\]
};

\node[subproc, below=of ind1] (ind2) {
\textbf{Valuar el portafolio en cada estado de la naturaleza}\\[-1mm]
Con tres bonos y una matriz con \(m\) estados:
\[
\omega=(s_1,s_2,s_3),\qquad
V_P(\omega)=V_{1,s_1}+V_{2,s_2}+V_{3,s_3}.
\]
Número de escenarios:
\[
m^3.
\]
Con 18 notches:
\[
18^3=5{,}832\ \text{distintos escenarios}.
\]
};

\node[subproc, below=of ind2] (ind3) {
\textbf{Obtener las probabilidades conjuntas de migración}\\[-1mm]
Si las migraciones fueran independientes, la distribución conjunta se obtiene multiplicando las probabilidades de transición:
\[
\begin{aligned}
p(\omega)
&=p_{A^-\to s_1}\,p_{BBB^-\to s_2}\\
&\quad\times p_{B^-\to s_3}
\end{aligned}
\]
};

\node[subproc, below=of ind3] (ind4) {
\textbf{Obtener el valor esperado y las desviaciones}\\[-1mm]
\[
\mu_P=\sum_{\omega}V_P(\omega)p(\omega).
\]
\[
\sigma_P=
\sqrt{\sum_{\omega}\left(V_P(\omega)-\mu_P\right)^2p(\omega)}.
\]
Pérdida en cada escenario:
\[
L(\omega)=V_{P,0}-V_P(\omega).
\]
};

\node[risk, below=of ind4] (ind5) {
Obtener el VaR al 99.9\%\\[-1mm]
Ordenar pérdidas \(L(\omega)\) y acumular probabilidades \(p(\omega)\).\\
\[
\begin{aligned}
VaR_{99.9\%}
&=\inf\{\ell:\Pr(L\leq \ell)\geq 0.999\}
\end{aligned}
\]
};

\node[output, below=of ind5] (ind6) {
Salida caso independiente\\
Distribución de pérdidas, pérdida esperada, desviación estándar y \(VaR_{99.9\%}\).\\
Interpretar como pérdida máxima al 99.9\% por cambios en la calidad crediticia.
};

% =========================================================
% RAMA DERECHA: DEPENDENCIA
% =========================================================

\node[process, below right=1.05cm and 5.1cm of tipo] (dep0) {
Caso II\\
Migraciones dependientes
};

\node[dep, below=of dep0] (dep1) {
\textbf{Modelar la correlación entre eventos de crédito}\\[-1mm]
La correlación de incumplimiento describe la tendencia de dos empresas a incumplir aproximadamente al mismo tiempo.\\
Obtener una matriz de correlaciones:
\[
R=
\begin{pmatrix}
1 & \rho_{12} & \rho_{13}\\
\rho_{21} & 1 & \rho_{23}\\
\rho_{31} & \rho_{32} & 1
\end{pmatrix}.
\]
La correlación de cópula puede aproximarse como la correlación entre los rendimientos de las acciones.
};

\node[dep, below=of dep1] (dep2) {
\textbf{Usar el modelo de cópula gaussiana}\\[-1mm]
Transformar cada tiempo al incumplimiento \(t_i\) en:
\[
x_i=N^{-1}\!\left(Q_i(t_i)\right),
\]
donde \(Q_i\) es la distribución acumulada de \(t_i\) y \(N^{-1}\) es la inversa de la distribución normal acumulada.
};

\node[dep, below=of dep2] (dep3) {
\textbf{Muestrear normales multivariadas}\\[-1mm]
Generar:
\[
X=(x_1,x_2,x_3)^\top\sim N(0,R).
\]
Por construcción, los \(x_i\) se distribuyen normal con media cero y desviación estándar unitaria, pero quedan correlacionados por \(R\).
};

\node[dep, below=of dep3] (dep4) {
\textbf{Convertir \(x_i\) a estado de crédito}\\[-1mm]
Usar los umbrales acumulados de la matriz de transición. Para cada bono \(i\):
\[
s_i=
\text{estado tal que }
\sum_{k< s_i}p_{i,k}
<
N(x_i)
\leq
\sum_{k\leq s_i}p_{i,k}.
\]
Así se obtiene un escenario dependiente:
\[
\omega=(s_1,s_2,s_3).
\]
};

\node[dep, below=of dep4] (dep5) {
\textbf{Repetir la valuación del portafolio}\\[-1mm]
Para cada escenario simulado:
\[
\begin{aligned}
V_P(\omega)&=V_{1,s_1}+V_{2,s_2}+V_{3,s_3}\\
L(\omega)&=V_{P,0}-V_P(\omega)
\end{aligned}
\]
Construir la distribución empírica de pérdidas con dependencia.
};

\node[risk, below=of dep5] (dep6) {
Obtener el VaR al 99.9\%\\[-1mm]
Ordenar pérdidas simuladas:
\[
L_{(1)}\leq L_{(2)}\leq \cdots \leq L_{(M)}.
\]
\[
VaR_{99.9\%}=L_{(\lceil0.999M\rceil)}.
\]
};

\node[output, below=of dep6] (dep7) {
Salida caso dependiente\\
Distribución de pérdidas con correlación, pérdida esperada, desviación estándar y \(VaR_{99.9\%}\).\\
Interpretar considerando que el riesgo de crédito no se puede diversificar por completo.
};

% =========================================================
% BLOQUE INFERIOR: n BONOS
% =========================================================

\coordinate (generalanchor) at ($(ind6.south)!0.5!(dep7.south)$);

\node[process, below=1.35cm of generalanchor] (general0) {
Procedimiento para un portafolio de \(n\) bonos
};

\node[subproc, below=of general0] (general1) {
\textbf{Con independencia}\\[-1mm]
1. Valuar cada bono en cada estado de crédito.\\
2. Formar todos los escenarios:
\[
\omega=(s_1,\ldots,s_n),\qquad \#\Omega=m^n.
\]
3. Obtener probabilidades conjuntas:
\[
p(\omega)=\prod_{i=1}^{n}p_{i,s_i}.
\]
4. Obtener \(V_P(\omega)\), \(\mu_P\), \(\sigma_P\) y el percentil 99.9\% de la distribución de pérdidas.
};

\node[subproc, right=1.0cm of general1] (general2) {
\textbf{Sin independencia}\\[-1mm]
1. Estimar cómo migran conjuntamente los obligados.\\
2. Usar una matriz de correlaciones \(R=(\rho_{ij})\).\\
3. Simular:
\[
X=(x_1,\ldots,x_n)^\top\sim N(0,R).
\]
4. Transformar percentil a percentil:
\[
x_i=N^{-1}(Q_i(t_i)).
\]
5. Convertir cada \(x_i\) a estado de crédito, valuar el portafolio y obtener el \(VaR_{99.9\%}\).
};

\node[note, below=of general1, xshift=3.15cm] (nota) {
Recordatorio operativo: CreditMetrics mide el riesgo de crédito por cambios de valor causados por migraciones en la calidad crediticia. 
Incorpora eventos de incumplimiento, mejoras y deterioros en la calidad crediticia; permite medir el beneficio de la diversificación o la potencial concentración de riesgo en la cartera.
};

\node[startstop, below=of nota] (fin) {
Fin\\
Reporte: \(VaR_{99.9\%}\) con independencia, \(VaR_{99.9\%}\) con dependencia e interpretación
};

\begin{pgfonlayer}{bg}
% =========================================================
% FLECHAS
% =========================================================
\draw[arrow] (inicio) -- (datos);
\draw[arrow] (datos) -- (insumos);
\draw[arrow] (insumos) -- (tipo);

\draw[arrow] (tipo) -| (ind0);
\draw[arrow] (ind0) -- (ind1);
\draw[arrow] (ind1) -- (ind2);
\draw[arrow] (ind2) -- (ind3);
\draw[arrow] (ind3) -- (ind4);
\draw[arrow] (ind4) -- (ind5);
\draw[arrow] (ind5) -- (ind6);

\draw[arrow] (tipo) -| (dep0);
\draw[arrow] (dep0) -- (dep1);
\draw[arrow] (dep1) -- (dep2);
\draw[arrow] (dep2) -- (dep3);
\draw[arrow] (dep3) -- (dep4);
\draw[arrow] (dep4) -- (dep5);
\draw[arrow] (dep5) -- (dep6);
\draw[arrow] (dep6) -- (dep7);

\draw[arrow] (ind6) |- (general0);
\draw[arrow] (dep7) |- (general0);
\draw[arrow] (general0) -- (general1);
\draw[arrow] (general0) -| (general2);
\draw[arrow] (general1) -- (nota);
\draw[arrow] (general2) |- (nota);
\draw[arrow] (nota) -- (fin);
\end{pgfonlayer}

\end{tikzpicture}
\end{adjustbox}
}

\vspace*{\fill}

\end{document}
